%%Template made by Uday Khankhoje for examinations using the exam template
%%Refer to the documentation http://www-math.mit.edu/~psh/exam/examdoc.pdf 
%%for lot more bells and whistles to the standard template shown below
\documentclass[a4paper,11pt,addpoints]{exam}
\usepackage[left=1.5cm,right=1.5cm,top=1.5cm,bottom=2cm]{geometry}
%\usepackage{mathrsfs}
\usepackage{graphicx,color}
\usepackage{epic,eepic}
%\usepackage{mathpazo}
\usepackage{url}
\usepackage{amsmath, amsthm, amssymb}
\pointsinmargin
\boxedpoints 
\renewcommand*\half{.5}
\usepackage{setspace}
\DeclareMathOperator{\vecc}{vec}
%\renewcommand{\vec}[1]{\ensuremath{\mathbf{#1}}}

\begin{document}
\noindent 
%%PART 1 of header
\begin{center}
	\vspace*{-3em}
\def\arraystretch{2.0}
\begin{tabular}{|p{0.7\linewidth}|p{0.2\linewidth}|}
\hline 
\textbf{EE 2025 Engineering Electromagnetics: End Semester Exam} & Marks obtained $\downarrow$ \\
\hline 
Date: 30.03.2020,\hspace{1cm}  Total questions: \textbf{\numquestions} \hspace{1cm} Total points: \textbf{\numpoints} &   \\ 
\hline 
\multicolumn{2}{|l|}{Roll No: \hspace{0.3\linewidth} Name: \hspace{0.3\linewidth} Time: 3 hrs} \\
\hline
\end{tabular} 
\end{center}
%%PART 2 of header, if you have too many questions, this may be a problem
%%if so, use \multirowgradetable{n}[questions], where n is the number of rows you want
%%or, switch to \gradetable[h][pages] instead, 
\begin{center}
\gradetable[h][questions]
\end{center}
%%PART 3 of header
\textbf{Instructions:
\begin{enumerate}
    \item Explain all steps clearly
    \item Attach this paper to the answer sheet
\end{enumerate}
}
%%toggle comment on next line to show/hide the answers
\printanswers
%%Now the actual paper!
\begin{questions}
\question[2] 
A first question for 2 points (indicated by the box on the left). If the solution environment has an optional spacing command given, then in the question-paper-only mode that much blank space will be left, like so.

\begin{solution}[1in]
This is the solution
\end{solution}

\question[1]
In this question we indicate no space in the solution environment, hence no space is left int he question-paper-only mode. Comment the ``{\tt{$\backslash$printanswers}}" line to show the solutions.

\begin{solution}
This is a solution
\end{solution}

\question
This question shows how to use parts with separate markings.
\begin{parts}
\part[3] This is the first part
\part[5] This is the second part
\end{parts}

\begin{solution}
\begin{enumerate}
    \item Here are the solutions to parts 1
    \item and 2.
\end{enumerate}

\end{solution}

\end{questions}
\end{document}